\section{Risk Neutral Asset Pricing}
\subsection{The One-Period Model}
\subsubsection{Introduction and Setup}
The market consists of a stock and a bank account over a single period $n\in
\left\{ 0,1 \right\}$. The bank account pays interests discretely at a rate
$r$. The stock price is a stochastic process ${\left( S_{n} \right)}_{n=0,1}$,
where $S_0\in \mathbb{R}^{+}_{*}$ and $S_1$ is a Bernoulli random variable such that:
\[S_1=S_0Z\quad\text{ where }\quad Z:\Omega \mapsto \left\{ u,d \right\},P[Z=u]=p_u, P[Z=d]=p_d.\]
For positive constants $d,u,p_u,p_d$, such that $d<u$, we refer to the market
model as the one-period binomial model.
\begin{definition}[Trading Strategy]
  We refer to a vector $h:=\left( x_0,y_0 \right)\in \mathbb{R}^{2}$ as
  the \textbf{portfolio} or \textbf{trading strategy}, where $x_0$ denotes the
  number of units of money in the bank account and $y_0$ denotes the
  number of units of stocks at time $t=0$ or period $n=0$.
\end{definition}
\begin{definition}[Value Process]
  We define the \textbf{value process} of a portfolio $h$ to be
  \[V_{n}^{h}:=x_0B_{n}+y_0S_{n},\qquad n=0,1.\]
\end{definition}
\begin{definition}[Arbitrage]
  A portfolio $h$ in the one-period binomial model is an \textbf{arbitrage portfolio} if the following
  conditions hold:
  \begin{itemize}
    \item $V_0^{h}=0$ 
    \item $P[V_1^{h}\ge 0]=1$
    \item $P[V_1^{h}> 0]>0$
  \end{itemize}
  That is, we are guaranteed a nonzero profit without risk.
\end{definition}
If no such portfolio exists, we say that the market is \textit{free of arbitrage}.
\begin{prop}[Free of Arbitrage]
  The one-period binomial model is free of arbitrage if and only if
  \[d\le 1+r\le u.\]
\end{prop}
\subsubsection{Risk-Neutral Measure}
With a market free of arbitrage in the one-period binomial model, we have $d\le 1+r\le u$ and
\[\exists q_u,q_d\in [0,1]\qquad \text{s.t.}\qquad 1+r=q_u\cdot u+q_d\cdot d\quad \text{and}\quad q_u+q_d=1.\] 
These conditions have the unique solution:
\[q_u=\frac{(1+r)-d}{u-d}\qquad\text{and}\qquad q_d=\frac{u-(1+r)}{u-d}.\]
We interpret $q_u$ and $q_d$ as a new probability measure $\mathbb{Q}$ where
the $\mathbb{Q}[Z=u]=q_u$ and $\mathbb{Q}[Z=d]=q_d$, of which the expected
value of the value of the stock is equal to that of the riskless bank account.
It follows that the value of the stock today is the discounted expected value
of the stock in the future.
This $\mathbb{Q}$ is known as the \textbf{risk-neutral measure}, (or
\textbf{martingale measure}), as we are indifferent between the bank account or
the stock under this measure. 
\begin{prop}[Existence of Risk-Neutral Measure]
  There exists such a measure if and only if the model is arbitrage free.
\end{prop}
\subsubsection{Contingent Claims}
With this market model, we are able to introduce derivative products and find the
fair prices for such objects.
\begin{definition}[Contingent Claim]
  A \textbf{contingent claim} is a stochastic variable $X$ of the form
  $X=\Phi(S_1)$ where $S_1$ is the stochastic variable representing the value
  of the stock.
\end{definition}
The real valued function $\Phi$ is known as the contract function.
\begin{notation}[Price Process of a Claim]
  We denote $\Pi\left( n;\Phi\left( S_1 \right)  \right) $ the price operator
  associated with the claim which returns the price of the claim at period $n$.
\end{notation}
\begin{definition}[Reachability of a Claim]
  We say that a contingent claim $X$ is \textbf{reachable} if there exits a portfolio
  $h$ such that
  \[P(V_1^{h}=X)=1.\]
  $h$ is known as the \textit{hedging portfolio} or a \textit{replicating
  portfolio}.
\end{definition}
By the law of one price, the hedging portfolio has the same value as the contingent claim:
\[\Pi\left( n;\Phi \right) =V_{n}^{h}.\]
By solving the equation at $n=1$ in the one-period binomial model,
\[V_1^{h}=\Phi=\Pi\left( 1;\Phi \right) \iff
  \begin{cases}x(1+r)+yS_0u=\Phi(S_0u)&Z=u\\x(1+r)+yS_0d=\Phi(S_0d)&Z=d\end{cases}\] 
which has a solution due to the no arbitrage condition $d<u$, $h^{*}=\left( x^{*},y^{*} \right) $ where:
\begin{align*}
  x^{*}=&\frac{1}{1+r}\frac{u\Phi\left( S_0d \right) -d\Phi\left( S_0u \right) }{u-d}\\
  y^{*}=&\frac{1}{S_0}\frac{\Phi\left( S_0u \right) -\Phi\left( S_0d \right) }{u-d}
\end{align*}
We have the price of the claim:
\[\Pi\left( 0;\Phi \right)=\frac{1}{1+r}E_{\mathbb{Q}}\left[  \Phi\left( S_1(Z) \right) \right].\] 
\begin{theorem}[Fair Price of a Contingent Claim]
  If the binomial model is free of arbitrage, then any claim $\Phi$ is replicable and
  the no-arbitrage price of $\Phi$ is given by
  \[\Pi\left( 0;\Phi \right)=\frac{1}{1+r}E_{\mathbb{Q}}\left[  \Phi\left( S_1\right) \right].\]
  where $\mathbb{Q}$ is the risk-neutral probability measure.
\end{theorem}
\begin{remark}
  The continuously compounded version of the one-period binomial measure
  follows, where $(1+r)$ is replaced by the continuously compounded interest
  $e^{rt}$
\end{remark}
\subsubsection{Risk-Free Portfolio Pricing}
We may also derive a portfolio such that there is no risk of the value of the
portfolio at its maturity. The value of the portfolio is independent of the
value of the underlying asset.\\
We find such a portfolio by solving for a portfolio $h=(0,\triangle)$ with a
short position in a European call with contract function $\Phi$ such that:
\[\triangle\cdot S_1(u)-\Phi(S_1(u))=\triangle\cdot S_1(d)-\Phi(S_1(d)).\]
\begin{remark}
  A riskless portfolio pays the exact same as a deposit in the bank account.
\end{remark}

\subsection{Multi-Period Model}
We extend the one-period binomial model into a multi-period binomial model, a
discrete-time model with $N \in \mathbb{N}$ periods and a running time index
$n=0,\ldots,N$. The time periods have a fixed time length $\triangle t$ and the
total time of the $N$ periods is $T=N\triangle t$.
\subsubsection{Setup}
The two underlying assets, namely a bond with price process
${(B_{n})}_{n=0,\ldots,N}$ and a stock with price process
${(S_{n})}_{n=0,\ldots,N}$, are defined as follows:
\[\text{Bank/Bond}=\begin{cases}B_0=1\\B_{n+1}=(1+r)B_{n}={(1+r)}^{n+1}\end{cases}\] 
\[\text{Bank/Bond}=\begin{cases}S_0=1\\S_{n+1}=S_{n}Z_{n}\end{cases}\] 
where we have a collection of i.i.d.\ Bernoulli r.v. $Z_0,\ldots,Z_{N-1}$.
\begin{definition}[Portfolio Strategy]
  We extend the portfolio strategy to a stochastic process
  \[\left\{ h_{n}=(x_{n},y_{n}) : n=1,\ldots,N\right\}.\]
  where each $h_{n}$ is a function of the stock prices. $h_0=h_1$ by
  convention. The corresponding \textit{value/wealth process} is
  \[V_{n}^{h}=x_{n}B_{n}+y_{n}S_{n}=x_{n}{(1+r)}^{n}+y_{n}S_{n}.\]
\end{definition}
\begin{definition}[Self-Financing Strategy]
  We say a portfolio strategy $h_{n}$ is a \textbf{self-financing strategy} if
  \[x_{n}(1+r)+y_{n}S_{n}=x_{n+1}(1+r)+y_{n+1}S_{n}.\]
  There is no extra capital into the portfolio at any time. We assume all
  portfolio strategies are self-financing.
\end{definition}
\begin{theorem}[Arbitrage-Free Market Model]
  For some $d<u$, the model is arbitrage-free if and only if $d\le 1+r\le u$,
  and there exists a unique risk-neutral measure $\mathbb{Q}$ whose defining probabilities
  $q_d$ and $q_u$ are defined by
  \[s=\frac{1}{1+r}E_{\mathbb{Q}}\left[ S_{n+1}|S_{n}=s \right].\] 
\end{theorem}
\subsubsection{Contingent Claims: European Claims}
\begin{notation}[Price Process of a Claim]
  We write, for a claim $\Phi$ written on $S$, for any $n\in \left\{ 0,\ldots,N \right\}$
  \[\left.\Pi\left( n;\Phi \right)\right|_{S_{n}}.\]
  the price of the option $\Phi$ at period $n$ when the stock takes the value
  $S_{n}$.
\end{notation}
\begin{prop}[Fair Price of a Contingent Claim]
  For a arbitrage-free $N$-period binomial market model and an European type claim $\Phi(S_N)$,
  \[\Pi\left( 0;\Phi\left( S_N \right)  \right) ={\left( e^{-r\triangle t} \right)}^{N}\sum_{k=0}^{N} \binom{N}{k}q_u^{k}q_d^{N-k}\Phi\left( S_0u^{k}d^{N-k} \right).\] 
\end{prop}

\subsubsection{Contingent Claims: American Options}
We consider claims of which we may exercise at any time. Exercising the
American option before maturity time is known as an \textit{early-exercise}.
\begin{definition}[Intrinsic Value]
  We define the intrinsic value of an American option at period $n$ as the
  market value (the payoff) of the option if one exercises the option at period
  $n$.
\end{definition}
If the intrinsic value of the American option is higher than the expected fair
price if it were an European option, one would exercise the option at that tree
node. Hence, we have:
\begin{itemize}
  \item At period $N$, the price of the option is exactly its pay-off.\\
  \item At periods $n<N$, the price of the option is the maximum between its
    expected future gain given risk neutrality, and its intrinsic value from an
    early exercise.
\[\left.\Pi\left( n;\Phi \right)\right|_{S_{n}=s}:=\max \left\{ \Phi(s),e^{-r\triangle t}\left( q_u\left.\Pi\left( n+1;\Phi \right) \right|_{S_{n+1}=us} + q_d\left.\Pi\left( n+1;\Phi \right) \right|_{S_{n+1}=ds}\right) \right\}.\] 
\end{itemize}
If the intrinsic value at period $n$ is less than the future expected value of
the option, the investor can expect to make a profit by waiting and will not
commit an early-exercise. If there is no early-exercise of the American call,
the price of the American call is the price as if it were a European call.
\begin{remark}
  The put-call parity does not hold for American options. We have, instead:
  \[S_0-K\le c_0-p_0\le S_0-Ke^{-rT}.\] 
\end{remark}

\subsection{Volatility and Model Calibration}
How should we pick $u,d$ such that it accurately models a stock? This depends
on the volatility of the stock, denoted $\sigma$, as a measure of uncertainty
of the returns of the stock.
\begin{remark}
  Stocks typically have a volatility betwen $15\%$ and $60\%$
\end{remark}
\subsubsection{Black-Scholes Model}
The returns of a stock are typically modelled over time as
\[\frac{\triangle S_t}{S_t}=\frac{S_{t+\triangle t}-S_t}{S_t}=\mu\triangle t+\sigma\sqrt{\triangle t}\varepsilon,\qquad \varepsilon\sim\mathcal{N}\left( 0,1 \right)\]
with $\mu$ representing the drift or trend. Additionally, the uncertainty about
a future stock price scales with the square root of ``how far into the
future''. We then have, in the binomial tree model:
\[S(t_{k+1})=S(t_k)\left( 1+\mu\triangle t+\sigma \sqrt{\triangle t}\varepsilon \right).\] 
\begin{remark}
  As a measure of uncertainty of the future stock price movements, higher
  volatility intuitively implies a greater chance of doing very well or very
  poorly. The two outcomes tend to offset each other as a investor of a stock.
  However, due to the limited downside risk of a option, investors of options
  benefit from greater volatility and \textit{the value of both calls and puts
  increases as volatility increases.}
\end{remark}
\subsubsection{Model Calibration}
Assuming we have the drift $\mu$ and volatility $\sigma$, how can we choose the
parameters of the model to accurately predict the market? We match the first
two moments of the Black-Scholes model and the Binomial model: Choosing $u,d,p_u,p_d$ such that:
\[E_\mathbb{P}\left[ S_1 \right]=E_\mathbb{P}\left[ S(t_1) \right]\qquad\text{and}\qquad \Var_\mathbb{P}\left[ S_1 \right]=\Var_\mathbb{P}\left[ S(t_1) \right].\] 
Matching the expectations, we have
\[S_0(u\cdot p_u+d\cdot\left( 1-p_u\right) )=S_0(1+e^{\mu\triangle t})\]
\[p_u=\frac{e^{\mu\triangle t}-d}{u-d}\]
Matching the variances, we have
\[S_0^{2}p_u\left( 1-p_u \right) {\left( u-d \right)}^{2}=\sigma^{2}S_0^{2}\triangle t\]
\[u=e^{\sigma \sqrt{\triangle t}}\qquad\text{and}\qquad d=e^{-\sigma \sqrt{\triangle t}}\]
\subsubsection{Risk-Neutral Measure}
The expectation under risk neutrality approaches the expectation under the
calibrated model as $r\to \mu$, but is not equal in general. The variance under
$\mathbb{Q}$ reproduces the calibrated model.\\
Moving between risk preferencees changes the expected growth rates, but the
volatilities remain the same (as $\triangle t\to0$). This is known as changing the (probability)
measure. [Girsanov's Theorem]

\subsection{Continuous-Time Multi-Period Binomial Model}
Since our model more accurately models the volatility as $\triangle t\to 0$, we
attempt to take the limits and compute the continuous-time case. 
\subsubsection{Setup}
Consider a market of continuous interest rate $r$ and volatility $\sigma$. Let
$\Phi\left( S(T) \right) $be a European Call or Put option with maturity $T$
written on a stock ${\left( S(t) \right) }_{t\in \left[ 0,T \right]}$. Consider $N$ sub-intervals
with time step $\triangle t=\frac{T}{N}$ and time points $t_{n}^{N}=n\times
\triangle t, n\in \left\{ 0,1,\ldots,N \right\}$.\\
We then have $u_N=e^{\sigma \sqrt{\triangle t}}, d_N=e^{-\sigma \sqrt{\triangle
t}}$ and probabilities
\[q_{u,N}=\frac{e^{r\triangle t}-d_N}{u-d}\qquad\text{and}\qquad q_{d,N}=1-q_{u,N}.\]
\begin{theorem}[Black-Scholes European Option Valuation Formula]
  The continuous time limit of the calibrated binomial for a call/put option of
  function $\Phi_c$/$\Phi_p$ with strike $K$ has fair prices:
  \[c_t=\lim_{\triangle t \to 0} \left.\Pi\left( 0,\Phi_c \right)
    \right|_{S_0}=S_0N\left( d_+ \right)-Ke^{-rT}N(d_{-}).\] 
  \[p_t=\lim_{\triangle t \to 0} \left.\Pi\left( 0,\Phi_p \right)
    \right|_{S_0}=Ke^{-rT}N(-d_{-})-S_0N\left( -d_+ \right).\] 
  where $d_+=\frac{\ln\left( \frac{S_0}{K}\right)+\left( r+\frac{\sigma^{2}}{2}
  \right) T  }{\sigma \sqrt{T}}$ and $d_-=d_+-\sigma \sqrt{T}=\frac{\ln\left(
  \frac{S_0}{K}\right)+\left( r-\frac{\sigma^{2}}{2} \right) T  }{\sigma
  \sqrt{T}}$ and $N(\cdot)$ the cumulative Normal distribution function.
\end{theorem}
